\documentclass{article}
\usepackage[utf8]{inputenc}
\usepackage[spanish]{babel}
\usepackage{tikz}
\usepackage{geometry}
\usepackage{adjustbox} % Indispensable para ajustar el tamaño

\geometry{a4paper, margin=0.5in, landscape}
\usetikzlibrary{shapes.multipart, shapes.geometric, positioning, fit, calc, arrows.meta, bending}

% --- DEFINICIÓN DEL ACTOR (Igual que antes) ---
\tikzset{
    actor/.pic={
        \node[circle, fill, minimum size=5mm] (head) {};
        \draw[thick] (head.south) -- ++(0,-1.5) coordinate (body);
        \draw[thick] (head.south)++(0,-0.5) -- ++(-0.5,0); 
        \draw[thick] (head.south)++(0,-0.5) -- ++(0.5,0);
        \draw[thick] (body) -- ++(-0.5,-0.8); 
        \draw[thick] (body) -- ++(0.5,-0.8);
        \node[below=1.2cm of head, align=center, font=\bfseries] {\tikzpictext};
    }
}

\begin{document}

\begin{figure}[h]
\centering

\begin{adjustbox}{max width=\textwidth, max height=0.9\textheight, keepaspectratio}

\begin{tikzpicture}[
    % Estilos globales
    node distance=2cm,
    usecase/.style={
        ellipse, 
        draw=black!80, 
        thick, 
        align=center, 
        fill=white, 
        minimum width=1cm, 
        minimum height=1.2cm,
        font=\small
    },
    usecase_complex/.style={
        ellipse split, 
        draw=black!80, 
        thick, 
        align=center, 
        fill=white, 
        minimum width=3.5cm,
        inner sep=3pt,
        font=\small
    },
    line/.style={-, thick, draw=gray!40!black},
    % Estilo para flechas de relación (dashed)
    dashed_arrow/.style={
        dashed, 
        -{Latex[open]}, 
        draw=black, 
        font=\scriptsize, 
        text=black
    }
]

% ==========================================
% 1. COLOCACIÓN DE NODOS (GRID ORDENADA)
% ==========================================

% --- ACTORES (Izquierda) ---
\pic [pic text=Comprador] (Buyer) at (-10, 6) {actor};
\pic [pic text=Vendedor] (Seller) at (-10, 0) {actor};
\pic [pic text=Administrador] (Admin) at (-10, -6) {actor};

% --- CASOS DE USO PRINCIPALES (Columna Central-Izquierda) ---
\node[usecase_complex] (Explorar) at (-3, 6) {
    \textbf{Explorar Catálogo}
    \nodepart{lower}
    \scriptsize points: filtro, detalle
};

\node[usecase] (Comprar) at (-3, 2.5) {Realizar Compra};

\node[usecase_complex] (GestionarProd) at (-3, 0) {
    \textbf{Gestionar Productos}
    \nodepart{lower}
    \scriptsize points: agregar, editar
};

\node[usecase] (VerPedidos) at (-3, -2.5) {Ver Pedidos};

\node[usecase_complex] (GestionarUser) at (-3, -6) {
    \textbf{Gestionar Usuarios}
    \nodepart{lower}
    \scriptsize points: aprobar, bloquear
};

% --- EXTENSIONES (Columna Derecha y Arriba) ---
% Usamos coordenadas relativas para alinear mejor

\node[usecase] (Filter) at (4, 7.5) {Filtrar Búsqueda};
\node[usecase] (Details) at (4, 6) {Ver Detalles};

\node[usecase] (Pay) at (4, 3.5) {Procesar Pago};

\node[usecase] (AddProd) at (4, 1.2) {Agregar Nuevo};
\node[usecase] (EditProd) at (4, -0.2) {Editar Existente};

\node[usecase] (UploadImg) at (10, 1.2) {Subir Imágenes}; % Mismo nivel que AddProd

\node[usecase] (Approve) at (4, -5) {Aprobar Vendedor};
\node[usecase] (Block) at (4, -9) {Bloquear Acceso};

% LOGIN (Ubicado estratégicamente abajo a la derecha pero accesible)
\node[usecase, fill=gray!10] (Login) at (15, -2) {Iniciar Sesión};


% ==========================================
% 2. CONEXIONES LIMPIAS
% ==========================================

% --- Conexiones Actores (Estilo Recto/Ortogonal) ---
% Comprador
\draw[line] (-10, 5.5) -- (Explorar.west);
\draw[line] (-10, 5.5) |- (Comprar.west); % |- significa bajar y luego horizontal

% Vendedor
\draw[line] (-10, -0.5) -- (GestionarProd.west);
\draw[line] (-10, -0.5) |- (VerPedidos.west);

% Admin
\draw[line] (-10, -6.5) -- (GestionarUser.west);


% --- Relaciones EXTEND (Curvas Suaves) ---
% Sintaxis: to [bend left=GRADOS] (destino)

% Explorar Extensions
\draw[dashed_arrow] (Filter) to [bend right=20] node[above, sloped] {$\ll$extend$\gg$} (Explorar);
\draw[dashed_arrow] (Details) -- node[below, sloped] {$\ll$extend$\gg$} (Explorar); % Recta porque están alineados

% Gestionar Productos Extensions
\draw[dashed_arrow] (AddProd) to [bend right=15] node[above, sloped] {$\ll$extend$\gg$} (GestionarProd);
\draw[dashed_arrow] (EditProd) -- node[below, sloped] {$\ll$extend$\gg$} (GestionarProd);

% Admin Extensions
\draw[dashed_arrow] (Approve) to [bend right=20] node[above, sloped] {$\ll$extend$\gg$} (GestionarUser);
\draw[dashed_arrow] (Block) to [bend left=20] node[below, sloped] {$\ll$extend$\gg$} (GestionarUser);


% --- Relaciones INCLUDE (Rutas complejas para evitar cruces) ---

% Comprar incluye Pago (directo)
\draw[dashed_arrow] (Comprar) -- node[above, sloped] {$\ll$include$\gg$} (Pay);

% AddProd incluye UploadImg (directo)
\draw[dashed_arrow] (AddProd) -- node[above, sloped] {$\ll$include$\gg$} (UploadImg);


% --- CONEXIONES A LOGIN (El mayor problema visual) ---
% Usamos curvas amplias para que vayan por fuera

% Comprar -> Login (Curva amplia por la derecha)
\draw[dashed_arrow] (Comprar) to [out=0, in=90] node[near end, sloped, above] {$\ll$include$\gg$} (Login.north);

% GestionarProd -> Login
\draw[dashed_arrow] (GestionarProd.south east) to [bend right=20] node[sloped, above] {$\ll$include$\gg$} (Login.west);

% VerPedidos -> Login
\draw[dashed_arrow] (VerPedidos) to [bend right=10] node[sloped, below] {$\ll$include$\gg$} (Login.west);

% GestionarUser -> Login (Curva amplia por abajo)
\draw[dashed_arrow] (GestionarUser.east) to [out=0, in=270] node[sloped, below] {$\ll$include$\gg$} (Login.south);


% --- LÍMITE DEL SISTEMA ---
\node[draw, thick, rectangle, rounded corners=1cm, dashed,
      fit=(Explorar) (GestionarUser) (UploadImg) (Login) (Filter) (Block),
      label={[anchor=south, font=\Large\bfseries, yshift=0.5cm]north:Sistema Silicon Trail}] (SystemBox) {};

\end{tikzpicture}

\end{adjustbox}
\caption{Diagrama de Casos de Uso }
\end{figure}

\end{document}